\scanpage{41}\mspage{40}
% L1–3
\uline{Characterization of an} \uline{ordered} \deleted{network-interval} \uline{interval}, \uline{in \deleted{so far} as a network}

% L4–5
\uline{by the datum} of
\diagram{\begin{tikzcd}[ampersand replacement=\&,row sep=0.75em,column sep=0.15em]
S \& {}\subset\unsure{\mathcal P|\mathscr P|\mathrm P}_2(L)\\
\makebox[0pt][c]{\(\operatorname{Drap}_2(L)\)}\arrow[u,equal] \&
\end{tikzcd}}
% L6–9
and for $\varepsilon=\{a,b\}\in S$, the decomposition of \deleted{$L$ into} $\complement_{\mathrm{mod}}\,\varepsilon=\{x\unsure{+|\in}L\mid\{\unsure{a;x|a',x|a,x}\},\{b,x\}\in S\}$ into ``connected components'' (by relation $R_\varepsilon$)

% L10–18
I began to do the exercise, but I have the impression that it is not worth the trouble of making it explicit to the end. It must have been done long ago (Hilbert~?) in the totally ordered \inserted{(and divisible)} case (in terms \unsure{at the very|so much at|at the very} least of the relation~: $s$ is between $a$ and $b$ \ldots\ )
